\begin{figure}[H]
\centering
\begin{tikzpicture}[x=3cm,y=3cm,font=\small,
    >=Stealth,
    initial/.style={blue!70!black},
    shifted/.style={orange!85!black},
    correction/.style={teal!75!black}]
\coordinate (reference) at (0,0);
\coordinate (initial) at (0,1.580);
\coordinate (shifted) at (-0.517,-0.104);
\coordinate (corner) at (-0.517,1.580);

\draw[step=0.5,gray!20,very thin] (-1,-0.5) grid (0.75,2);
\draw[->,gray!70!black] (-1,-0.5) -- (0.9,-0.5);
\draw[->,gray!70!black] (-1,-0.5) -- (-1,2.12);
\foreach \tick in {-1,-0.5,0,0.5}{
    \draw[gray!70!black] (\tick,-0.5) -- (\tick,-0.54)
        node[below,font=\footnotesize] {$\tick$};
}
\foreach \tick in {0,0.5,1,1.5,2}{
    \draw[gray!70!black] (-1,\tick) -- (-1.04,\tick)
        node[left,font=\footnotesize] {$\tick$};
}
\node[below] at (-0.08,-0.74) {$x-X_{10,1}$ (píxeles)};
\node[rotate=90] at (-1.42,0.75) {$y-Y_{10,1}$ (píxeles)};

\draw[correction,dashed,->] (initial) -- (corner)
    node[midway,above=5pt] {$\overline{\Delta x}=-0.517$};
\draw[correction,dashed,->] (corner) -- (shifted)
    node[midway,left=5pt,rotate=90,anchor=south]
    {$\overline{\Delta y}=-1.684$};
\draw[correction,very thick,->,shorten <=3pt,shorten >=4pt]
    (initial) -- (shifted);
\node[correction,anchor=west,align=left] at (0.08,0.85)
    {Traslación\\\texttt{imshift}};

\draw[gray!70!black,densely dotted,thick] (shifted) -- (reference);
\fill[initial] (initial) circle (2.6pt);
\node[initial,right=6pt] at (initial) {\texttt{im011}};
\fill[shifted] (shifted) circle (2.6pt);
\node[shifted,below left=3pt] at (shifted) {\texttt{s011}};
\draw[black,thick] (-0.055,0) -- (0.055,0);
\draw[black,thick] (0,-0.055) -- (0,0.055);
\node[above right=3pt] at (reference) {\texttt{im010}};
\end{tikzpicture}

\smallskip
{\small
\begin{tabular}{lrr}
\toprule
Posición de la estrella 1 & $X$ (píxeles) & $Y$ (píxeles) \\
\midrule
Referencia: \texttt{im010} & 289.986 & 83.967 \\
Entrada: \texttt{im011} & 289.986 & 85.547 \\
Salida predicha: \texttt{s011} & 289.469 & 83.863 \\
\bottomrule
\end{tabular}
}
\caption{Esquema geométrico de la corrección con \texttt{imshift},
aplicada a la estrella 1. El origen de las coordenadas relativas es su
posición en \texttt{im010}.}
\label{fig:imshift}
\end{figure}